%% notes-tex/wet-lab/hydrolysate-figure-gate/sections/4-state.tex
%% SOURCE: hand-authored state record. Measured values are macros of
%% tables/gate_numbers.tex. Job identifiers and dates are from
%% notes/experiments.040-inhibitor-synergy-wetlab.md (sections of 2026.10.10) and from
%% the directory listings of experiments/038-env-chemgen-vanacloig-cgt-corrected/results/
%% factorized/ at the time the sheets were drawn. Everything under "expect" is a
%% hypothesis and is labeled as one.

\clearpage
\section{State on the day the sheets were drawn\secstatus{todo}}
\label{sec:state}

\pfstep{What had finished} The processor round of this experiment: the four
pre-registered claims scored, which is every panel marked on file above. The ladder of
the corrected-store campaign, giving main \textbf{d} and SI-1 \textbf{a}. The learning
curve of the earlier campaign, giving main \textbf{f} in its ridge arm. The pooled table
and its axes, giving main \textbf{b}.

\pfstep{What was running} The second factorized round of the corrected-store campaign,
submitted in the small hours of the day the sheets were drawn on two GilaHyper cards: one
job for the bilinear control, one for the environment encoder. Each had written two folds
when the panel script read the directory, the control at a pooled median of
\hydbilinearMedian{} and the encoder at \hydenvencMedian{} over
\hydenvencCompounds{} held-out compounds. Those are two folds of a round, not a result,
and main \textbf{e} is marked partial for that reason. The sheet re-reads the directory
every time it is drawn, so the panel fills in as the round writes.

\pfstep{What had not started} The first mixture round, which is every panel of the bottom
row in its model arm. The trainer and its configuration groups are owned by the
implementation thread and are not part of this document.

\pfstep{What to expect, labeled} Hypothesis (untested): the environment encoder ties
nested ridge on the corrected store, as it did on the earlier store, so main \textbf{e}
closes with the two within a hundredth of each other and the figure's public-benchmark
row says that architecture is not where the gain is. Hypothesis (untested): the encoder's
compound-cold profiles for the two inhibitors in store read no better than ridge, so main
\textbf{i} keeps its shape and the predicted profiles stay generic. Hypothesis
(untested): a dose-aware host head fitted on the titration run beats the highest single
agent on growth, which reads \hydhsaAuroc{}, and does not beat Loewe, which reads
\hydloeweAuroc{}, because Loewe already uses the fitted single-agent curves and the model
has no information about this host that those curves do not carry. If that last one holds
the honest figure is the model-free one, and the bottom row says so in the caption rather
than in a footnote.

\pfstep{What would change the figure's claim} A model arm that beats
\hydloeweAuroc{} on growth, or \hydblissSpearman{} on fitness over the combinations that
grew, moves the bottom row from a validation of arithmetic to a result about transfer. A
model arm that does neither still earns the row, because the comparison is the point: the
rules are computed from six titrations on the same strain, and a model trained on
\hydnPoolCompounds{} public compounds that cannot beat them is the measurement the
section needs to report.
